\documentclass[11pt]{article}
\usepackage[T1]{fontenc}
\usepackage{lmodern}
\usepackage{microtype}
\usepackage[margin=1in]{geometry}
\usepackage{amsmath,amssymb,amsthm,mathtools}
\usepackage{graphicx}
\usepackage{xcolor}
\usepackage[colorlinks=true,linkcolor=blue!55!black,citecolor=blue!55!black,urlcolor=blue!55!black]{hyperref}

\newtheorem{theorem}{Theorem}
\newtheorem{corollary}[theorem]{Corollary}
\newtheorem{remark}[theorem]{Remark}
\newcommand{\essinf}{\operatorname*{ess\,inf}}
\newcommand{\Lop}{\mathbb L}
\newcommand{\Aop}{\mathcal A}

\title{A rank-one counterexample to the graphon Perron converse}
\author{Automated mathematics run \texttt{fa\_banach\_001}\\
Lane 10 --- candidate result for expert review}
\date{9 August 2026}

\begin{document}
\maketitle

\begin{center}
\fcolorbox{black!30}{gray!7}{\parbox{0.91\textwidth}{
\textbf{Status.} Candidate counterexample to the literal open problem in
Bonnet--Pouradier Duteil--Sigalotti, together with an exact characterization
inside the full-support rank-one class. The proof is elementary and appears
complete, but the interpretation caveat in Section~\ref{sec:caveat} is
important.
}}
\end{center}

\section{Source question}

Bonnet, Pouradier Duteil, and Sigalotti ask on printed page 25 of
arXiv:2111.03900, after Theorem 4.7, whether the finite-dimensional Perron
converse persists for graphon Laplacians \cite{source}. Their question is:

\begin{quote}
Is it true that an interaction topology is a countable disjoint union of
strongly connected components whenever there exists an element
$v\in L^\infty(I,\mathbb R_+^*)$ such that $\Lop^*v=0$?
\end{quote}

\begin{figure}[ht]
  \centering
  \includegraphics[width=\textwidth]{figures/open_problem_crop.png}
  \caption{The open problem in the source paper, printed page 25. The crop is
  rendered from the local copy \texttt{source\_paper.pdf}.}
\end{figure}

Here $I=[0,1]$, an adjacency kernel $a\in L^\infty(I\times I,[0,1])$
defines
\[
 (\Aop x)(i)=\int_I a(i,j)x(j)\,dj,\qquad
 d(i)=\int_I a(i,j)\,dj,\qquad
 \Lop=M_d-\Aop.
\]
The source's Definition 4.5 requires a strongly connected topology to satisfy
both qualitative reachability through row supports and the quantitative bound
$\inf_{i\in I}d(i)>0$.  A disjoint union must consist of at most countably many
closed positive-measure components and satisfy the corresponding normalized
degree bound with one uniform positive constant.

\section{Proof intuition}

The adjoint-null equation is a balance law, whereas Definition 4.5 also asks
for a uniform interaction rate.  Rank-one kernels separate these two demands
completely.  If $a(i,j)=p(i)q(j)$, the adjoint equation can be solved explicitly:
the stationary weight is proportional to $q/p$.  In contrast, the degree is
proportional to $p$.  Taking $p=q$ makes the stationary weight constant no
matter how close $p$ comes to zero.  Full support prevents a partition from
hiding the low-degree region inside different components.

\section{Rank-one characterization}

All assertions below are understood up to null sets, as is standard for
$L^\infty$ kernels.  For the final continuous counterexample, ordinary and
essential supports and infima give the same conclusion.

\begin{theorem}[Full-support rank-one class]\label{thm:rankone}
Let $p,q\in L^\infty(I)$ satisfy $0<p,q\leq1$ almost everywhere, and put
\[
 a(i,j)=p(i)q(j),\qquad Q=\int_I q(j)\,dj.
\]
Then:
\begin{enumerate}
\item A function $v\in L^\infty(I)$ with $v>0$ almost everywhere satisfies
      $\Lop^*v=0$ if and only if
      \[
      v=c\,\frac{q}{p}
      \]
      for some $c>0$. Consequently such a bounded strictly positive weight
      exists exactly when $q/p\in L^\infty(I)$; with $\int_Iv=1$ it is unique.
\item If $q>0$ almost everywhere, every nonzero row has support $I$. Hence a
      Definition 4.5 disjoint-union decomposition exists if and only if the
      whole space is its single component, which occurs if and only if
      \[
      \essinf_{i\in I}p(i)>0.
      \]
\end{enumerate}
In particular, bounded positive adjoint-null weights do not imply the
quantitative strong-connectivity conclusion.
\end{theorem}

\begin{proof}
For $x\in L^2(I)$,
\[
 (\Aop x)(i)=p(i)\int_Iq(j)x(j)\,dj,
 \qquad d(i)=Qp(i).
\]
The adjoint of the rank-one operator is
\[
 (\Aop^*y)(i)=q(i)\int_Ip(j)y(j)\,dj.
\]
Writing $C=\int_Ip(j)v(j)\,dj$, the equation $\Lop^*v=0$ is therefore
\[
 Qp(i)v(i)=q(i)C \quad\text{for almost every }i.
\]
Since $p,v>0$ almost everywhere, $C>0$, and this identity is equivalent to
$v=(C/Q)(q/p)$. Conversely, if $v=cq/p$, then
$C=c\int_Iq=cQ$, so the displayed equation holds. Normalizing the integral
fixes $c=(\int_Iq/p)^{-1}$. This proves the first assertion.

For almost every $i$, $p(i)>0$, so the row $a(i,\cdot)$ is a positive scalar
multiple of $q$. Since $q>0$ almost everywhere, its support is all of $I$.
The qualitative path condition in Definition 4.5(a) is automatic: any two
eligible points are joined in one step.

Suppose $\{I_n\}$ were a disjoint-union decomposition in the sense of that
definition. Each $I_n$ has positive measure. Choosing an eligible
$i\in I_n$ gives
\[
 I=\operatorname{supp}a(i,\cdot)\subset I_n
\]
up to null sets. Thus there can only be one component, namely $I$ itself.
On this component Definition 4.5(b) is precisely
\[
 0<\inf_{i\in I}d(i)=Q\inf_{i\in I}p(i),
\]
or, invariantly, $\essinf p>0$. The converse is immediate: when this bound
holds, $I$ itself satisfies both conditions of Definition 4.5.
\end{proof}

\section{The minimal counterexample}

\begin{corollary}\label{cor:counterexample}
The open problem has a negative answer as literally stated.  Define
\[
 a(i,j)=ij\qquad (i,j\in[0,1]).
\]
Then $v\equiv1$ belongs to $L^\infty(I,\mathbb R_+^*)$ and satisfies
$\Lop^*v=0$, but the topology is not a disjoint union of strongly connected
components in the sense of Definition 4.5.
\end{corollary}

\begin{proof}
This is Theorem~\ref{thm:rankone} with $p(i)=q(i)=i$ (the value at the null
point $i=0$ is immaterial). Directly,
\[
 d(i)=\int_0^1ij\,dj=\frac{i}{2},
 \qquad
 (\Aop1)(i)=\frac{i}{2},
\]
and $\Aop$ is self-adjoint because $a$ is symmetric. Hence
$\Lop^*1=\Lop1=d-\Aop1=0$. For each $i>0$, the continuous function
$j\mapsto ij$ has support $I$, so an admissible closed component containing
$i$ must be all of $I$. But
\[
 \inf_{i\in I}d(i)=\inf_{i\in I}\frac{i}{2}=0,
\]
contradicting Definition 4.5(b).
\end{proof}

\section{Interpretation caveat and corrected rank-one statement}
\label{sec:caveat}

The example is qualitatively strongly connected: for $i,j>0$, the edge weight
$a(i,j)$ is positive, and the exceptional point $i=0$ is negligible. It fails
only the uniform degree condition built into the source's terminology.
Therefore:
\begin{itemize}
\item Under Definition 4.5 exactly as printed, Corollary~\ref{cor:counterexample}
      is a full counterexample.
\item If ``strongly connected components'' in the open problem was intended to
      mean only qualitative mutual reachability, this example does not refute
      that weaker intended statement. It instead identifies the missing
      quantitative hypothesis.
\item In the full-support rank-one class, the exact repair is
      $\essinf p>0$. The stationary equation alone controls $q/p$ and cannot
      supply this bound.
\end{itemize}

There is also a same-source warning about novelty. In Section 5.4 the authors
give a different symmetric, qualitatively connected kernel whose degree
infimum is zero \cite{source}. Symmetry automatically yields $\Lop^*1=0$.
Thus the obstruction to the literal open problem is already latent in the
source paper, although that section does not state that it answers the open
problem. The rank-one calculation above isolates the mechanism and gives a
complete elementary family, but priority for the bare negative answer should
be treated as low.

\section{Verification and novelty bounds}

The proof is symbolic and uses no numerical computation. The following checks
were made explicitly:
\begin{enumerate}
\item the adjoint of the rank-one operator was computed from its defining
      $L^2$ pairing;
\item substituting $v=cq/p$ reproduces the same scalar $C=cQ$;
\item full row support rules out every multi-component partition allowed by
      Definition 4.5;
\item for $a(i,j)=ij$, both the literal and essential degree infima vanish.
\end{enumerate}

A bounded search on 9 August 2026 covered the run's registry, solution,
attempt, and proof-gap indexes; arXiv id 2111.03900; the exact open-problem
phrase; the paper title combined with ``Perron''; and web searches for positive
stationary graphon weights and strong components. It found the source and
later citing papers, but no separate paper explicitly resolving this converse.
This is not an exhaustive MathSciNet or zbMATH priority search. The latent
Section 5.4 example is the principal novelty limitation.

\paragraph{Human-review recommendation.}
Check first whether the authors intended Definition 4.5 in full in the open
problem. If so, the counterexample is immediate and the rank-one theorem is a
sharp diagnosis. If only qualitative reachability was intended, retain the
rank-one theorem as a statement correction and do not advertise the weaker
problem as solved.

\begin{thebibliography}{9}
\bibitem{source}
B. Bonnet, N. Pouradier Duteil, and M. Sigalotti,
\emph{Consensus Formation in First-Order Graphon Models with Time-Varying
Topologies}, Math. Models Methods Appl. Sci. 32 (2022), no. 11, 2121--2188;
arXiv:2111.03900; DOI: \href{https://doi.org/10.1142/S0218202522500518}{10.1142/S0218202522500518}.
\end{thebibliography}

\end{document}
